\documentclass[11pt]{article}
\usepackage[margin=1in]{geometry}
\usepackage{amsmath,amssymb,amsthm}
\usepackage{graphicx}
\usepackage[hidelinks]{hyperref}

\newtheorem{theorem}{Theorem}
\newtheorem{lemma}{Lemma}

\newcommand{\A}{\mathcal A}
\newcommand{\B}{\mathcal B}
\newcommand{\C}{\mathbb C}
\newcommand{\dist}{\operatorname{dist}}
\newcommand{\HH}{\mathcal H}
\newcommand{\K}{\mathcal K}
\newcommand{\N}{\mathbb N}
\newcommand{\Toe}{\mathcal T}

\title{A Toeplitz-Algebra Counterexample to a C*-Distance Formula Question}
\author{Run \texttt{fa\_banach\_001}, lane 11}
\date{2026-06-19}

\begin{document}
\maketitle

\noindent\textbf{Status.} \texttt{counterexample\_likely\_valid}.

\section{Source Question}

The source paper is Grover--Singla, \emph{Birkhoff-James orthogonality and
applications: A survey}, arXiv:2005.07399 \cite{GroverSingla}.
On pages 15--16, after recalling Rieffel's distance formula, the authors ask
whether the following formula holds. If $\A$ is a unital irreducible
$C^*$-algebra, $a\in\A$, and $\B$ is a unital $C^*$-subalgebra of $\A$, is it
true that
\begin{equation}\label{eq:question}
2\,\dist(a,\B)
=
\sup\{\|au-ua\|:\ u\in\A,\ u \text{ unitary, and } bu=ub
\text{ for all } b\in\B\}?
\end{equation}
The source states that it is open whether \eqref{eq:question} is true.

\begin{center}
\includegraphics[width=0.97\linewidth]{figures/open_problem_crop.png}
\end{center}

\section{Proof Intuition}

The right hand side of \eqref{eq:question} only sees unitaries in the relative
commutant $\B'\cap\A$. If $\B$ is already irreducible inside $\A$, this
relative commutant can be only the scalar multiples of the identity. Then the
right hand side is forced to be zero for every $a$. A proper irreducible
subalgebra therefore gives a counterexample as soon as one chooses
$a\notin\B$.

The standard Toeplitz algebra provides a very concrete example. It is generated
by the unilateral shift $S$ on $\ell_2(\N)$, has scalar commutant, and is a
proper subalgebra of $B(\ell_2)$. A simple diagonal unitary
$D e_n=(-1)^n e_n$ lies outside it because Toeplitz-algebra elements have
compact commutator with $S$, whereas $[D,S]$ is noncompact.

\section{Counterexample}

\begin{theorem}
Formula \eqref{eq:question} is false.
\end{theorem}

\begin{proof}
Let $\HH=\ell_2(\N)$ with orthonormal basis $(e_n)_{n\geq 0}$, and let
$\A=B(\HH)$. Then $\A$ is a unital irreducible $C^*$-algebra. Let
$S\in B(\HH)$ be the unilateral shift, $Se_n=e_{n+1}$, and set
\[
  \B=C^*(S)\subset B(\HH).
\]
This is the unital Toeplitz algebra. Define a diagonal unitary $D$ by
\[
  D e_n = (-1)^n e_n \qquad (n\geq 0).
\]

We first compute the relative commutant. If $T\in B(\HH)$ commutes with every
element of $\B$, then $T$ commutes with both $S$ and $S^*$. Since
$S^*e_0=0$, we have
\[
  S^*T e_0 = T S^* e_0 =0,
\]
so $T e_0\in\ker S^*=\C e_0$. Write $T e_0=\lambda e_0$. Because $T$ commutes
with $S$,
\[
  T e_n = T S^n e_0 = S^n T e_0 = \lambda e_n
\]
for every $n$. Hence $T=\lambda I$. Therefore
\[
  \B'\cap \A = \C I.
\]
The only unitary elements of $\A$ commuting with all of $\B$ are thus
$\lambda I$, $|\lambda|=1$. For such a unitary,
\[
  D(\lambda I)-(\lambda I)D=0.
\]
Consequently the right hand side of \eqref{eq:question}, with $a=D$, is $0$.

It remains to show that the left hand side is positive. We prove $D\notin\B$.
Let
\[
  \mathcal E=\{X\in B(\HH): XS-SX \text{ is compact}\}.
\]
This is a norm-closed unital algebra: if $X,Y\in\mathcal E$, then
\[
  (XY)S-S(XY)=X(YS-SY)+(XS-SX)Y,
\]
which is compact, and norm-closedness follows because the compact operators
are norm closed. Also $S\in\mathcal E$, and
\[
  S^*S-SS^* = I-SS^*
\]
is the rank-one projection onto $\C e_0$, so $S^*\in\mathcal E$. Since
$C^*(S)$ is the norm-closed unital algebra generated by $S$ and $S^*$, it
follows that
\[
  C^*(S)\subseteq \mathcal E.
\]

But
\[
  (DS-SD)e_n = D e_{n+1}-S((-1)^n e_n)
  = (-1)^{n+1}e_{n+1}-(-1)^n e_{n+1}
  = -2(-1)^n e_{n+1}.
\]
Thus $DS-SD$ sends the orthonormal basis $(e_n)$ to vectors of constant norm
$2$. It is not compact. Hence $D\notin\mathcal E$, and therefore
$D\notin C^*(S)=\B$.

Since $\B$ is norm closed and $D\notin\B$, we have $\dist(D,\B)>0$. Thus the
left hand side of \eqref{eq:question} is $2\dist(D,\B)>0$, while the right
hand side is $0$. This contradicts \eqref{eq:question}.
\end{proof}

\section{Verification Notes}

The proof is analytic and does not rely on computation. The included script
\texttt{code/toeplitz\_sanity\_check.py} verifies the finite-section shadow of
the construction: the commutator of the alternating diagonal with the shift has
stable norm $2$, matching the noncompactness mechanism in the proof.

Novelty checks were bounded as follows. The run indexes were searched for
\texttt{2005.07399}, \texttt{Birkhoff-James}, \texttt{distance formula},
\texttt{C*-subalgebra}, and \texttt{Toeplitz}. Web searches on 2026-06-19 for
exact and close variants of \eqref{eq:question} found the source arXiv paper
but no later paper explicitly resolving the question.

\section{Limitations and Review Recommendation}

This counterexample disproves formula \eqref{eq:question} as stated when
``unital irreducible $C^*$-algebra'' is interpreted in the standard represented
sense, allowing $\A=B(\ell_2)$. It does not address possible repaired formulas
with additional hypotheses, for example requiring a sufficiently large
relative commutant or replacing unitaries in $\A\cap\B'$ by unitaries in a
cyclic representation as in Rieffel's theorem.

Human review should focus on the intended meaning of ``irreducible'' in the
source question and on the elementary Toeplitz-commutant/properness checks.

\begin{thebibliography}{9}
\bibitem{GroverSingla}
P. Grover and S. Singla,
\emph{Birkhoff-James orthogonality and applications: A survey},
arXiv:2005.07399, 2020.
\end{thebibliography}

\end{document}
